\documentclass[12pt]{article}

% ======================================================
% Codificación e idioma
% ======================================================

\usepackage[utf8]{inputenc}
\usepackage[T1]{fontenc}
\usepackage{multicol}
\usepackage[spanish,es-nodecimaldot]{babel}

% =====================================================
% Matemáticas
% =====================================================

\usepackage{amsmath,amssymb,amsfonts}

% =====================================================
% Formato y utilidades
% =====================================================

\usepackage[letterpaper,margin=2.2cm]{geometry}
\usepackage{enumitem}
\usepackage{array}
\usepackage{tabularx}
\usepackage{xcolor}
\usepackage{fancyhdr}
\usepackage{lastpage}
\usepackage{hyperref}

% =====================================================
% Hipervínculos
% =====================================================

\hypersetup{
  colorlinks=true,
  linkcolor=blue,
  urlcolor=blue,
  citecolor=blue
}

% =====================================================
% Formato de párrafos
% =====================================================

\setlength{\parindent}{0pt}
\setlength{\parskip}{0.05em}

% =====================================================
% Datos editables del examen
% =====================================================

\newcommand{\institucion}{Tecnológico Nacional de México}
\newcommand{\materia}{Cálculo Integral}
\newcommand{\nombreexamen}{Primer Examen Parcial}
\newcommand{\unidad}{Unidad 1: Integral Indefinica}
\newcommand{\docente}{Carlos Ernesto Martínez}
\newcommand{\fechaexamen}{}
\newcommand{\duracion}{60 minutos}
\newcommand{\puntajetotal}{100 puntos}

% =====================================================
% Encabezado y pie de página
% =====================================================

\setlength{\headheight}{10pt}

\pagestyle{fancy}
\fancyhf{}

\fancyhead[L]{\institucion}
\fancyhead[R]{\materia}

\fancyfoot[L]{\docente}
\fancyfoot[C]{Página \thepage\ de \pageref{LastPage}}
\fancyfoot[R]{\nombreexamen}

\renewcommand{\headrulewidth}{0.1pt}
\renewcommand{\footrulewidth}{0.1pt}

% =====================================================
% Contador y ambiente para problemas
% =====================================================

\newcounter{problema}

\newenvironment{problema}[1]
{
  \refstepcounter{problema}
  \par\medskip
  \noindent
  \textbf{Problema \theproblema.}
  \textbf{[#1 puntos]}
  \par\smallskip
}
{
  \par\medskip
}

% =====================================================
% Comando para espacio de respuesta
% =====================================================

\newcommand{\espaciorespuesta}[1]{
  \vspace{#1}
}

% =====================================================
% Documento
% =====================================================

\begin{document}

% =====================================================
% Encabezado principal
% =====================================================

\begin{center}
    {\Large\textbf{\institucion}}\\[0.1em]
    {\large\textbf{\materia}}\\[0.1em]
    {\large\textbf{\nombreexamen}}\\[0.1em]
    {\normalsize\unidad}
\end{center}

\vspace{0.15em}

% =====================================================
% Datos del estudiante
% =====================================================

\begin{tabularx}{\textwidth}{>{\bfseries}l X >{\bfseries}l X}
Nombre: & \hrulefill &
No. de control: & \hrulefill \\[1.2em]

Grupo: & \hrulefill &
Fecha: & \hrulefill \\[1.2em]

Calificación: & \hrulefill &

\end{tabularx}



% =====================================================
% Información general
% =====================================================

% =====================================================
% Instrucciones
% =====================================================

\subsection*{Instrucciones}
Muestre de manera clara y ordenada todo el procedimiento. Justifique sus resultados cuando se solicite. Las respuestas sin procedimiento no serán consideradas par la calificaci\'on. Sí se permite uso de calculadora, mostrar al inicio del examen.

% =====================================================
% REACTIVOS
% =====================================================


\begin{problema}{10}*
Identifique claramente la región, las curvas que la delimitan y el intervalo de interés. Dibuje la región limitada por $y=9-x^2$, el eje $x$ y las rectas $x=0$ y $x=3$.
\end{problema}

\begin{problema}{10}
Para $f(x)=\sqrt{x+1}$ en $[0,3]$, calcule aproximaciones izquierda y derecha con $n=6$.
\end{problema}


\begin{problema}{15}
En los siguientes problemas use $n$ rectángulos y obtenga una fórmula para
$A_n$. Después calcule $A=\lim_{n\to\infty}A_n.$ Muestre explícitamente $\Delta x$, $x_i$, $f(x_i)$ y la suma correspondiente  $y=2x+3$, $1\le x\le4$.
\end{problema}


\begin{problema}{10}
Determine el valor exacto de la integral definida interpretándola como área de una región plana $\displaystyle \int_0^2\sqrt{4-x^2}\,dx$.

\end{problema}


\begin{problema}{10}
Determine, con aproximación de centésimos, un valor de $c$ que satisfaga
$\int_a^b f(x)\,dx=f(c)(b-a).$ $\displaystyle \int_1^4(x^2+4x+5)\,dx$.
\end{problema}

\begin{problema}{10}
*Determine $F'(x)$.  
\begin{enumerate}
\begin{multicols}{2}
\item $\displaystyle F(x)=\int_2^x\sqrt{1+t^2}\,dt$.
\item $F(x)=\int_{-2}^x e^t\,dt$
\end{multicols}
\end{enumerate}
\end{problema}

\begin{problema}{20}
*Evalúe las siguientes integrales. \textbf{Resolver a lo m\'as dos ejercicios}.

\begin{enumerate}
\begin{multicols}{2}
\item $\displaystyle \int_{-1}^3(3x^2+5x-1)\,dx$.
\item $\displaystyle \int_0^1\frac{y^2+2y}{\sqrt[3]{y^3+3y^2+4}}\,dy$.
\item $\displaystyle \int_0^1\sin(\pi x)\cos(\pi x)\,dx$.
\item $\displaystyle \int_0^1\frac{y^2+2y}{\sqrt[3]{y^3+3y^2+4}}\,dy$.
\item $\displaystyle \int_0^3(x+2)\sqrt{x+1}\,dx$.

\end{multicols}
\end{enumerate}

\end{problema}

\begin{problema}{20}
Calcule mostrando la antiderivada y la evaluación en los límites, \textbf{Resolver a lo m\'as dos ejercicios}.

\begin{enumerate}
\begin{multicols}{2}
\item Evalúe $\int_1^4\frac{1}{\sqrt{x}}\,dx$.
\item Evalúe $\int_0^{\pi/2}(\sin x+\cos x)\,dx$.
\item Evalúe $\int_1^4\left(x^2+\sqrt{x}+x^{-2}\right)\,dx$.
\end{multicols}
\end{enumerate}

\end{problema}

\begin{problema}{15}
Obtenga la derivada. \textbf{Resolver a lo m\'as dos ejercicios}.

\begin{enumerate}
\begin{multicols}{3}

\item $\frac{d}{dx}\int_2^x\frac{1}{t^4+4}\,dt.$

\item $\frac{d}{dx}\int_{-x}^{x}\cos(t^2+1)\,dt.$

\item $\frac{d}{dx}\int_1^{x^3}\sqrt[3]{t^2+1}\,dt.$
\end{multicols}
\end{enumerate}

\end{problema}

\begin{problema}{15}
*\textbf{Resolver a lo m\'as dos ejercicios}. Calcule
\begin{enumerate}
\begin{multicols}{3}
\item $\sum_{k=1}^{6}(2k+1)^2$.
\item $\sum_{k=1}^{20}3k(k^{2}+2)^2.$
\item $\sum_{k=1}^{n}\left(2^{k}-2^{k-1}\right)$
\end{multicols}
\end{enumerate}

\end{problema}




%\end{multicols}
\espaciorespuesta{6cm}


\end{document}